% ============================================================
% example-paper.tex — harryopo-paper.cls 单栏模式完整示例
% 展示：中文标题/摘要、多作者、算法伪代码、代码高亮、基金信息
% ============================================================
\documentclass{harryopo-paper}

% --- 论文标题 ---
\title{基于注意力机制的轻量化图像分类方法研究}

% --- 作者信息 ---
\author{张三\textsuperscript{1} \quad 李四\textsuperscript{2} \quad 王五\textsuperscript{1}\\%
       {\small\textsuperscript{1} 计算机学院，深圳 518172}\\
       {\small\textsuperscript{2} 南方科技大学 计算机科学与工程系，深圳 518055}}
% --- 日期 ---
\date{2026年6月21日}

% --- 基金信息 ---
\fundproject{基金项目：国家自然科学基金（No. 12345678）；深圳市科技计划项目（No. JCYJ20240001）}

% --- 作者简介 ---
\authorbio{张三（1990--），男，博士，副教授，主要研究方向为计算机视觉、深度学习。E-mail: zhangsan@harryopo.edu.cn}

\begin{document}

% ==== 标题 + 中文摘要 ====
\maketitle

\begin{abstract}
图像分类是计算机视觉领域的基础任务之一。近年来，深度学习技术的快速发展为图像分类带来了革命性的突破。然而，高性能模型通常伴随着巨大的参数量和计算成本，难以部署到资源受限的边缘设备。本文提出了一种基于注意力机制的轻量化图像分类方法，在保持较高准确率的同时显著降低了模型复杂度。实验结果表明，所提方法在ImageNet数据集上的Top-1准确率达到78.3\%，参数量仅为3.2M，FLOPs为0.8G，优于现有的MobileNetV3和ShuffleNetV2等轻量化模型。本文还对注意力机制的作用原理进行了深入分析，并通过消融实验验证了各模块的有效性。
\end{abstract}

\keywords{图像分类；注意力机制；轻量化网络；深度学习}

% ==== 正文 ====
\section{引言}

图像分类是计算机视觉中最基本的任务之一，其目标是将输入图像分配到预定义的类别中\cite{lecun2015deep}。传统方法依赖手工设计的特征提取器，泛化能力有限。深度学习技术的兴起——特别是卷积神经网络（CNN）——彻底改变了这一局面。

自AlexNet\cite{krizhevsky2012imagenet}在2012年ImageNet挑战赛中取得突破以来，研究者们提出了大量CNN架构。ResNet\cite{he2016deep}通过残差连接解决了深层网络训练困难的问题，EfficientNet\cite{tan2019efficientnet}则通过复合缩放策略实现了效率与精度的平衡。

然而，这些高性能模型往往包含数百万甚至数千万参数，难以在移动设备、嵌入式系统等资源受限场景中部署。为此，研究者们提出了多种轻量化设计策略。

本文的主要贡献如下：
\begin{enumerate}
    \item 提出了一种融合通道注意力和空间注意力的轻量化模块，在保持低计算量的同时有效增强了特征表达
    \item 设计了渐进式通道压缩策略，逐步减少特征通道数，显著降低参数量
    \item 在ImageNet和CIFAR-100数据集上进行了充分的实验验证
\end{enumerate}

\section{相关工作}

\subsection{轻量化卷积神经网络}

轻量化网络设计是当前的研究热点之一。MobileNet系列\cite{howard2017mobilenets}使用深度可分离卷积替代标准卷积，将计算量降低约8--9倍。ShuffleNet\cite{zhang2018shufflenet}引入通道混洗操作，在分组卷积后增强通道间的信息流动。EfficientNet\cite{tan2019efficientnet}提出了复合缩放策略。

\subsection{注意力机制}

注意力机制最早在自然语言处理领域取得了巨大成功，随后被引入计算机视觉领域。SE-Net通过全局平均池化和全连接层学习通道间的依赖关系。CBAM同时使用通道注意力和空间注意力来增强特征表示。

\subhead{通道注意力}通道注意力关注"什么"是有意义的特征通道，通过全局池化聚合空间信息后学习通道维度的权重。

\subhead{空间注意力}空间注意力关注"哪里"是关键区域，通过通道维度的池化后学习空间位置的权重分布。

\section{本文方法}

\subsection{整体架构}

本文提出的轻量化网络架构如图\ref{fig:architecture}所示。网络由三个主要阶段组成：初始特征提取阶段、深度特征学习阶段和分类预测阶段。

\begin{figure}[htbp]
    \centering
    \colorbox{gray!10}{\parbox{0.85\textwidth}{\centering\vspace{2cm}%
        {\large\color{gray} [网络架构示意图]}\\[0.5cm]
        \small 输入 → Stem → Stage1 → Stage2 → Stage3 → Pooling → FC → 输出
    \vspace{2cm}}}
    \caption{本文提出的轻量化图像分类网络架构}
    \label{fig:architecture}
\end{figure}

\subsection{轻量化注意力模块}

本文设计的轻量化注意力模块（LAM）同时融合了通道注意力和空间注意力，其结构如算法\ref{alg:lam}所示。

\begin{algorithm}[htbp]
\caption{轻量化注意力模块前向传播}
\label{alg:lam}
\KwIn{输入特征图 $X \in \mathbb{R}^{C \times H \times W}$}
\KwOut{增强后的特征图 $Y \in \mathbb{R}^{C \times H \times W}$}

% 通道注意力分支
$g_c \leftarrow \text{GAP}(X)$ \Comment{全局平均池化 → $C \times 1 \times 1$}

$g_c \leftarrow \text{Conv1d}_{k=3}(\text{ReLU}(\text{Conv1d}_{k=3}(g_c)))$ \Comment{一维卷积通道交互}

$A_c \leftarrow \sigma(g_c)$ \Comment{Sigmoid 激活}

$X_c \leftarrow A_c \odot X$ \Comment{通道维加权}

% 空间注意力分支
$g_s \leftarrow \text{Mean}(X_c, \text{dim}=1)$ \Comment{通道维取均值 → $1 \times H \times W$}

$g_s \leftarrow \text{Conv}_{7\times7}([g_s; \text{Max}(X_c, \text{dim}=1)])$ \Comment{拼接均值与最大值}

$A_s \leftarrow \sigma(g_s)$ \Comment{Sigmoid 激活}

$Y \leftarrow A_s \odot X_c$ \Comment{空间维加权}

\Return{$Y$}
\end{algorithm}

\subhead{通道注意力分支}通道注意力分支首先通过全局平均池化将空间信息压缩为一个向量，然后使用一维卷积在相邻通道间进行信息交互，最后使用 Sigmoid 函数生成通道权重。

\subhead{空间注意力分支}空间注意力分支在通道维度上分别计算均值和最大值，拼接后使用大卷积核（$7\times7$）捕获空间上下文信息，生成空间注意力图。

\subsection{渐进式通道压缩}

传统轻量化网络通常在不同阶段使用固定的通道数配置。本文提出了渐进式通道压缩策略，在网络深层逐步减少通道数，同时保持足够的特征维度。

设第$i$个阶段的输出通道数为$C_i$，则：
\begin{equation}
C_i = C_0 \cdot \alpha^{i} \quad (i = 1, 2, \dots, N)
\label{eq:channel}
\end{equation}
其中$C_0$为基础通道数，$\alpha \in (0, 1]$为压缩因子，$N$为阶段总数。本文中取$C_0=64$，$\alpha=0.75$。

\section{实验结果}

\subsection{实验设置}

所有实验在PyTorch框架下完成，使用NVIDIA RTX 3090 GPU训练。优化器为SGD with momentum，初始学习率0.1，cosine衰减策略，batch size为256，共训练300个epoch。

\subsection{ImageNet分类结果}

表\ref{tab:comparison}展示了本文方法与主流轻量化模型在ImageNet数据集上的性能对比。

\begin{table}[htbp]
\centering
\begin{tabular}{lcccc}
\toprule
\textbf{方法} & \textbf{参数量(M)} & \textbf{FLOPs(G)} & \textbf{Top-1(\%)} & \textbf{Top-5(\%)} \\
\midrule
MobileNetV2 & 3.5 & 0.30 & 72.0 & 91.0 \\
MobileNetV3-Large & 5.4 & 0.22 & 75.2 & 92.2 \\
ShuffleNetV2 & 2.3 & 0.15 & 72.6 & 90.6 \\
EfficientNet-B0 & 5.3 & 0.39 & 77.1 & 93.3 \\
GhostNet & 5.2 & 0.14 & 73.9 & 91.4 \\
\midrule
\textbf{本文方法} & \textbf{3.2} & \textbf{0.18} & \textbf{78.3} & \textbf{94.0} \\
\bottomrule
\end{tabular}
\caption{ImageNet数据集上的性能对比}
\label{tab:comparison}
\end{table}

从表\ref{tab:comparison}可以看出，本文方法以3.2M的参数量和0.18G的FLOPs，取得了78.3\%的Top-1准确率，在精度和效率之间达到了最优平衡。

\subsection{消融实验}

表\ref{tab:ablation}展示了各模块的消融实验结果。

\begin{table}[htbp]
\centering
\begin{tabular}{ccccc}
\toprule
\textbf{Baseline} & \textbf{LAM} & \textbf{PCC} & \textbf{Params(M)} & \textbf{Top-1(\%)} \\
\midrule
\checkmark &  &  & 3.0 & 76.1 \\
\checkmark & \checkmark &  & 3.4 & 77.5 \\
\checkmark &  & \checkmark & 2.8 & 76.8 \\
\checkmark & \checkmark & \checkmark & 3.2 & \textbf{78.3} \\
\bottomrule
\end{tabular}
\caption{消融实验结果}
\label{tab:ablation}
\end{table}

注：LAM为轻量化注意力模块，PCC为渐进式通道压缩策略。实验结果表明，两种策略均可独立提升模型性能，且具有协同效应。

\section{结论}

本文提出了一种基于注意力机制的轻量化图像分类方法。通过在MobileNetV3的基础上引入融合通道和空间注意力的轻量化模块，并辅以渐进式通道压缩策略，在保持低计算量的同时显著提升了分类精度。未来的工作将探索将该方法扩展到目标检测和语义分割等下游任务。

% ==== 基金信息 ====
\printfunding

% ==== 参考文献 ====
\begin{thebibliography}{9}
\bibitem{lecun2015deep} LeCun Y, Bengio Y, Hinton G. Deep learning[J]. Nature, 2015, 521(7553): 436-444.
\bibitem{krizhevsky2012imagenet} Krizhevsky A, Sutskever I, Hinton G E. ImageNet classification with deep convolutional neural networks[C]. NeurIPS, 2012.
\bibitem{simonyan2014very} Simonyan K, Zisserman A. Very deep convolutional networks for large-scale image recognition[J]. arXiv:1409.1556, 2014.
\bibitem{he2016deep} He K, Zhang X, Ren S, et al. Deep residual learning for image recognition[C]. CVPR, 2016.
\bibitem{huang2017densely} Huang G, Liu Z, Van Der Maaten L, et al. Densely connected convolutional networks[C]. CVPR, 2017.
\bibitem{howard2017mobilenets} Howard A G, Zhu M, Chen B, et al. MobileNets: Efficient convolutional neural networks for mobile vision applications[J]. arXiv:1704.04861, 2017.
\bibitem{zhang2018shufflenet} Zhang X, Zhou X, Lin M, et al. ShuffleNet: An extremely efficient convolutional neural network for mobile devices[C]. CVPR, 2018.
\bibitem{tan2019efficientnet} Tan M, Le Q V. EfficientNet: Rethinking model scaling for convolutional neural networks[C]. ICML, 2019.
\end{thebibliography}

% ==== 作者简介 ====
\printauthorbio

\end{document}
